% Source: Weinberg GR, physical PDF pp. 185--187
%         (printed pp. 161--163).
% Coverage: Section 7.4, Coordinate Conditions; equations (7.4.1)--(7.4.9).
% Figures/tables/footnotes: none.
% Status: source-reviewed and compile-clean.
% Uncertainties: none.

\subsection{Coordinate Conditions}\label{sec:7.4}

The symmetric tensor \(G_{\mu\nu}\) has 10 independent components, so
Einstein's field equations \eqref{eq:7.1.13} comprise 10 algebraically
independent equations. The unknown metric tensor also has 10 algebraically
independent components, and at first sight one would think that the Einstein
equations (with appropriate boundary conditions) would suffice to determine
the \(g_{\mu\nu}\) uniquely. However, this is not so. Although algebraically
independent, the 10 \(G_{\mu\nu}\) are related by four differential
identities, the Bianchi identities (see Eq.~\eqref{eq:6.8.3}):
\[
  \nabla_\mu G^{\mu\nu}=0.
\]
Thus there are not 10 functionally independent equations, but only
\(10-4=6\), leaving us with four degrees of freedom in the 10 unknowns
\(g_{\mu\nu}\). These degrees of freedom correspond to the fact that if
\(g_{\mu\nu}\) is a solution of Einstein's equation, then so is
\(g'_{\mu\nu}\), where \(g'_{\mu\nu}\) is determined from \(g_{\mu\nu}\) by
a general coordinate transformation \(x\to x'\). Such a coordinate
transformation involves four arbitrary functions \(x'^\mu(x)\), giving to
the solutions of Eq.~\eqref{eq:7.1.13} just four degrees of freedom.

The failure of Einstein's equations to determine \(g_{\mu\nu}\) uniquely is
closely analogous with the failure of Maxwell's equation to determine the
vector potential \(A_\mu\) uniquely. When written in terms of the vector
potential, Maxwell's equations read
\begin{equation}
  \Box A_\mu-\partial_\mu\partial_\nu A^\nu=-J_\mu.
  \tag{7.4.1}\label{eq:7.4.1}
\end{equation}
(See Eqs.~\eqref{eq:2.7.6} and \eqref{eq:2.7.11}.) There are four equations
for the four unknowns, but they do not determine \(A_\mu\) uniquely, because
the left-hand sides of these equations are related by a differential identity
analogous to the Bianchi identities:
\[
  \partial^\mu
  \left(
    \Box A_\mu-\partial_\mu\partial_\nu A^\nu
  \right)
  \equiv0.
\]
Thus the number of functionally independent equations is really only
\(4-1=3\), and there is one degree of freedom in the solution for the four
\(A_\mu\). This degree of freedom of course corresponds to gauge invariance;
given any solution \(A_\mu\), we can find another solution
\(A'_\mu=A_\mu+\partial_\mu\Lambda\), with \(\Lambda\) arbitrary.

The ambiguity in the solutions of Maxwell's and Einstein's equations can be
removed by main force. In the case of Maxwell's equation we do this by
choosing a particular gauge. For instance, given any solution \(A_\mu\), we
can always construct a solution \(A'_\mu\) such that
\begin{equation}
  \partial_\mu A'^\mu=0
  \tag{7.4.2}\label{eq:7.4.2}
\end{equation}
by setting
\[
  A'_\mu=A_\mu+\partial_\mu\Phi,
\]
where \(\Phi\) is defined by
\[
  \Box\Phi=-\partial_\mu A^\mu.
\]
Such a solution is said to be in the Lorenz gauge. The condition
\eqref{eq:7.4.2}, when added to the three independent equations
\eqref{eq:7.4.1}, completes a system of four equations that, with appropriate
boundary conditions, will generally determine the four \(A_\mu\) uniquely.
In the same way, we can eliminate the ambiguity in the metric tensor by
adopting some particular coordinate system. The choice of a coordinate system
can be expressed in four \emph{coordinate conditions}, which, when added to
the six independent Einstein equations, determine an unambiguous solution.

One particularly convenient choice of a coordinate system is represented by
the \emph{harmonic coordinate conditions}
\begin{equation}
  \Gamma^\lambda
  \equiv
  g^{\mu\nu}\Gamma^\lambda{}_{\mu\nu}
  =
  0.
  \tag{7.4.3}\label{eq:7.4.3}
\end{equation}
To see that it is always possible to choose a coordinate system in which this
holds, we recall the transformation equations of the affine connection:
\[
  \Gamma'^\lambda{}_{\mu\nu}
  =
  \frac{\partial x'^\lambda}{\partial x^\rho}
  \frac{\partial x^\tau}{\partial x'^\mu}
  \frac{\partial x^\sigma}{\partial x'^\nu}
  \Gamma^\rho{}_{\tau\sigma}
  -
  \frac{\partial x^\rho}{\partial x'^\nu}
  \frac{\partial x^\sigma}{\partial x'^\mu}
  \frac{\partial^2x'^\lambda}
       {\partial x^\rho\partial x^\sigma}.
\]
(See Eq.~\eqref{eq:4.5.8}.) Contracting this with \(g'^{\mu\nu}\), we find
\begin{equation}
  \Gamma'^\lambda
  =
  \frac{\partial x'^\lambda}{\partial x^\rho}\Gamma^\rho
  -
  g^{\rho\sigma}
  \frac{\partial^2x'^\lambda}
       {\partial x^\rho\partial x^\sigma}.
  \tag{7.4.4}\label{eq:7.4.4}
\end{equation}
Hence if \(\Gamma^\rho\) does not vanish, we can always define a new
coordinate system \(x'^\lambda\) by solving the second-order partial
differential equations
\[
  g^{\rho\sigma}
  \frac{\partial^2x'^\lambda}
       {\partial x^\rho\partial x^\sigma}
  =
  \frac{\partial x'^\lambda}{\partial x^\rho}\Gamma^\rho,
\]
and Eq.~\eqref{eq:7.4.4} then gives \(\Gamma'^\lambda=0\) in the
\(x'\)-system.

The four conditions \eqref{eq:7.4.3} are of course not generally covariant,
since their purpose is to remove the ambiguity in the metric tensor owing to
the general covariance of the Einstein equations. Although we cannot write
them as covariant equations, they can be put in a somewhat more elegant form
by expressing the affine connection in terms of the metric tensor:
\[
  \Gamma^\lambda
  =
  \frac12g^{\mu\nu}g^{\lambda\kappa}
  \left(
    \partial_\nu g_{\kappa\mu}
    +\partial_\mu g_{\kappa\nu}
    -\partial_\kappa g_{\mu\nu}
  \right).
\]
We recall that
\[
  g^{\lambda\kappa}\partial_\nu g_{\kappa\mu}
  =
  -g_{\kappa\mu}\partial_\nu g^{\lambda\kappa},
  \qquad
  \frac12g^{\mu\nu}\partial_\kappa g_{\mu\nu}
  =
  \frac1{\sqrt{-g}}\partial_\kappa\sqrt{-g}.
\]
(See Eq.~\eqref{eq:4.7.5}.) This then gives
\begin{equation}
  \Gamma^\lambda
  =
  -\frac1{\sqrt{-g}}
  \partial_\kappa\left(\sqrt{-g}\,g^{\lambda\kappa}\right),
  \tag{7.4.5}\label{eq:7.4.5}
\end{equation}
and the harmonic coordinate conditions read
\begin{equation}
  \partial_\kappa\left(\sqrt{-g}\,g^{\lambda\kappa}\right)=0.
  \tag{7.4.6}\label{eq:7.4.6}
\end{equation}

We are now in a position to explain the term ``harmonic coordinates.'' A
function \(\phi\) is said to be harmonic if \(\Box\phi\) vanishes, where
\(\Box\) is the invariant d'Alembertian, defined by
\begin{equation}
  \Box\phi
  \equiv
  \nabla_\kappa\left(g^{\lambda\kappa}\nabla_\lambda\phi\right).
  \tag{7.4.7}\label{eq:7.4.7}
\end{equation}
Using Eqs.~\eqref{eq:4.7.1}, \eqref{eq:4.7.7}, and
\eqref{eq:7.4.5}, this is
\begin{equation}
  \Box\phi
  =
  g^{\lambda\kappa}\partial_\lambda\partial_\kappa\phi
  -\Gamma^\lambda\partial_\lambda\phi.
  \tag{7.4.8}\label{eq:7.4.8}
\end{equation}
If \(\Gamma^\lambda=0\), then the coordinates are themselves harmonic
functions,
\begin{equation}
  \Box x^\mu=0,
  \tag{7.4.9}\label{eq:7.4.9}
\end{equation}
thus justifying our application of the adjective ``harmonic'' to such
coordinate systems.

In the absence of gravitational fields, the obvious harmonic coordinate
system is that of Minkowski, in which \(g^{\lambda\kappa}=\eta^{\lambda\kappa}\)
and \(g=-1\), so that Eq.~\eqref{eq:7.4.6} is satisfied trivially. In the
presence of weak gravitational fields the harmonic coordinate systems may be
pictured as nearly Minkowskian. Another related advantage of the harmonic
coordinate condition is that, as shown in Chapters~9 and~10, its use produces
a very great simplification in the weak-field equations, similar to the
simplification brought to Maxwell's equations by use of the Lorenz gauge.
